\section{Methodology}
This section condenses five notebook-trained components; Appendix~\ref{app:extended} expands pedagogical detail teammates can lift verbatim into the final write-up.

\subsection{Text featurization}
Each catalogue row forms a \emph{text blob} by repeating tokens from tags, genres, categories, and a truncated description chunk; fields are up-weighted (tags highest) to emphasise game-play semantics over prose \cite{deerwester1990lsi}. A scikit-learn \texttt{TfidfVectorizer} with character n-grams on word tokens, English stop-words, sublinear TF, and document-frequency pruning yields a sparse term-document matrix \(X\in\mathbb{R}^{N\times V}\) \cite{sklearn2011}.

\subsection{Latent semantic factors (LSA)}
Truncated singular value decomposition with \(k=128\) components approximates
\begin{equation}
X \approx U_k \Sigma_k V_k^\top ,
\end{equation}
where rows of the projected matrix serve as dense embeddings. L2-normalization enforces cosine geometry comparable to collaborative-filtering cosine similarities on item factors. These vectors power nearest-neighbour scoring, downstream classifiers, regressors, and clustering.

\subsection{Hybrid recommender}
For anchor item \(i\) and candidate \(j\), we compute TF--IDF cosine similarity \(s^{\text{tf}}_{ij}\), latent cosine \(s^{\text{lsa}}_{ij}=\langle \ell_i,\ell_j\rangle\), and normalized popularity \(p_j\in[0,1]\). The blend
\begin{equation}
h_{ij}= w_{\text{lsa}} s^{\text{lsa}}_{ij} + w_{\text{text}} s^{\text{tf}}_{ij} + w_{\text{pop}} p_j
\end{equation}
defaults to richer latent weighting for item--item mode and slightly higher popularity prior for cold-start user queries. Optional quantile floors suppress ultra-niche neighbours below a popularity percentile. A \emph{content richness} multiplier penalises tag-sparse rows in profile mode, mitigating degenerate short queries.

\subsection{Explainability}
Post-hoc explanations enumerate shared genres and rank shared tags by inverse document frequency, surfacing why two titles scored similarly without SHAP overhead adequate for demos.

\subsection{Player-type weak supervision}
Eight archetypes (Explorer, Narrative Nerd, Tinkerer, Trophy Hunter, Thrill-Seeker, Grinder, Speedrunner, Competitor) are encoded as deterministic labelling functions over tags/genres/categories. The resulting binary label matrix is noisy but cheap. A one-vs-rest logistic regression on LSA vectors predicts membership probabilities; class weighting counters imbalance.

\subsection{Popularity regression without leakage}
The target is Bayesian-shrunk \texttt{popularity}. Features include LSA factors, scaled release year, counts of tags/genres/categories/platforms, description length, and curated boolean indicators (multiplayer flags, indie genre, Steam achievement categories, early access, VR). Explicitly excluded are raw review counts, thumbs tallies, and Metacritic---they would trivialise the task.

An XGBoost histogram regressor with early stopping serves as the primary model \cite{chen2016xgboost}. Residual analysis highlights cult classics (under-predicted) versus bait-and-switch tag spam (over-predicted).

\subsection{Hidden genres via KMeans}
KMeans partitions LSA row vectors into interpretable groups; silhouette scores on subsampled data select \(k\). Cluster identities emerge from tag \emph{lift}: the ratio of intra-cluster tag frequency to global frequency. Auto-names stitch the top high-lift tags subject to count thresholds to avoid naming clusters after single-digit noise tokens.

\subsection{Release-month seasonality}
Parsing release timestamps yields month-of-year histograms per theme (Horror via tags, Sports/Racing/Strategy/RPG via genres). We chart a \emph{seasonal index}: monthly mean divided by theme-wide mean. Forecast experiments compare a flat \(1/12\) month-share baseline against historical-shape averages and report mean absolute error in percentage points on withheld years.
